% 01-description.tex - Section 1: Description

\chapter{Description}
\label{sec:description}

\section{Purpose}

This document describes the interface control for the Next-Generation Digital Processor (NDP) subsystem of the Long Wavelength Array at the New Mexico site (LWA-NA). The NDP is responsible for digitizing, channelizing, and processing the analog signals from the Front End Electronics (FEE) via the Analog Signal Processor (ASP), producing beamformed and correlated data products.

\section{Scope}

This Interface Control Document (ICD) defines:
\begin{itemize}
    \item Physical interfaces to and from NDP
    \item Monitor and control (M\&C) interfaces
    \item Output data interfaces and formats
    \item Safety interfaces
    \item Command and response message formats
    \item Management Information Base (MIB) entries
\end{itemize}

\section{System Overview}

The NDP subsystem performs the following functions:
\begin{itemize}
    \item Analog-to-digital conversion at 196~MHz sample rate
    \item F-engine processing: Polyphase filter bank (PFB) channelization with 4096 channels
    \item X-engine processing: Power line interference flagging (Appendix~\ref{app:powerline})
    \item X-engine processing: Beamforming (up to 4 dual-polarization beams)
    \item X-engine processing: Correlation (full visibility matrix)
    \item X-engine processing: Streaming complex spectra output from all antennas (TBS)
    \item X-engine processing: Triggered complex spectra dumps from all antennas (TBT)
    \item T-engine processing: Conversion from frequency domain to time domain for DRX output
\end{itemize}

\subsection{Hardware Architecture}

The NDP system consists of three main hardware components:

\subsubsection{ZCU102 F-engine Boards}

The F-engines perform analog-to-digital conversion and frequency channelization using Xilinx ZCU102 FPGA boards.  Table~\ref{tab:fengine-config} summarizes the F-engine configuration.

\begin{table}[htbp]
    \centering
    \caption{F-engine Configuration}
    \label{tab:fengine-config}
    \begin{tabular}{ll}
        \toprule
        \textbf{Parameter} & \textbf{Value} \\
        \midrule
        Number of boards & 16 \\
        Inputs per board & 32 (16 stands $\times$ 2 polarizations) \\
        Sample rate & 196~MHz \\
        ADC resolution & 8 bits \\
        FFT channels & 4096 \\
        Channel bandwidth & 23.926~kHz \\
        \bottomrule
    \end{tabular}
\end{table}

\subsubsection{GPU Servers (X-engines)}

The X-engines perform beamforming and correlation using GPU-accelerated servers running the Bifrost framework.  Table~\ref{tab:xengine-config} lists the X-engine parameters.

\begin{table}[htbp]
    \centering
    \caption{X-engine Configuration}
    \label{tab:xengine-config}
    \begin{tabular}{ll}
        \toprule
        \textbf{Parameter} & \textbf{Value} \\
        \midrule
        Number of servers & 4 \\
        Pipelines per server & 4 \\
        Total pipelines & 16 \\
        Processing framework & Bifrost \\
        \bottomrule
    \end{tabular}
\end{table}

\subsubsection{T-engines}

The T-engines convert frequency-domain beamformed data back to the time domain for DRX-format output, producing Mark~5C compatible packets.  Table~\ref{tab:tengine-config} lists the T-engine parameters.

\begin{table}[htbp]
    \centering
    \caption{T-engine Configuration}
    \label{tab:tengine-config}
    \begin{tabular}{ll}
        \toprule
        \textbf{Parameter} & \textbf{Value} \\
        \midrule
        Number of servers & 1 \\
        Pipelines per server & 4 (one per beam) \\
        Total pipelines & 4 \\
        Processing framework & Bifrost \\
        \bottomrule
    \end{tabular}
\end{table}

\subsection{System Capacity}

Table~\ref{tab:system-capacity} summarizes the overall system capacity.

\begin{table}[htbp]
    \centering
    \caption{NDP System Capacity}
    \label{tab:system-capacity}
    \begin{tabular}{ll}
        \toprule
        \textbf{Parameter} & \textbf{Value} \\
        \midrule
        Maximum stands & 256 \\
        Polarizations per stand & 2 \\
        Total inputs & 512 \\
        Maximum beams & 4 (dual-polarization) \\
        Frequency channels & 4096 \\
        Maximum tunings per beam & 2 \\
        \bottomrule
    \end{tabular}
\end{table}

\section{Observing Modes}

\subsection{DRX (Beam Recording)}

The primary observing mode provides up to 4 dual-polarization beams with configurable tuning frequencies. Each beam can have up to 2 independent frequency tunings. The DRX output is in Mark~5C compatible format.

\subsection{COR (Correlator)}

The correlator mode cross-correlates all 512 inputs to produce a full visibility matrix. The integration time is configurable via the \cmd{COR} command.

\subsection{TBS (Transient Buffer -- Streaming)}

Continuous streaming mode that outputs frequency-domain spectra for all 512 inputs. TBS is designed for 100\% duty cycle operation with a narrow bandwidth selection ($\sim$200~kHz at filter code 8). This mode can run concurrently with other observing modes.

\subsection{TBT (Transient Buffer -- Triggered)}

Triggered dump mode that captures frequency-domain spectra with wider bandwidth than TBS. TBT trades duty cycle for increased bandwidth, making it suitable for capturing transient events with more spectral coverage.

\section{Document Organization}

This document is organized as follows:
\begin{itemize}
    \item \textbf{Section~\ref{sec:description}} (this section): General description and scope
    \item \textbf{Section~\ref{sec:abbreviations}}: Abbreviations and definitions
    \item \textbf{Section~\ref{sec:physical}}: Physical interface specifications
    \item \textbf{Section~\ref{sec:monitor-control}}: Monitor and control interface
    \item \textbf{Section~\ref{sec:output}}: Output data interfaces
    \item \textbf{Section~\ref{sec:safety}}: Safety interface
    \item \textbf{Section~\ref{sec:references}}: References
\end{itemize}
